\honest{The Simple Truth}{Eliezer Yudkowsky}{2008}

``This essay is meant to restore a naive view of truth.'' When a seller of snake oil answers a clinical study by asking ``what do you mean by `true'?'', you do not need a rigorous definition to keep the word. People fell off cliffs before anyone had the equations of gravity. The essay is not ``an encyclopedic reference''; I want people arguing on mailing lists to read it and get back to their subject.

Then a story. I am a shepherd before recorded history, and if I stop searching early I lose a sheep about one night in ten. Divination sticks and clairvoyance do not help. So I drop a pebble into a bucket for each sheep that goes out, and take one out for each sheep that comes in. When the bucket is empty, the pasture is empty, provided the bucket starts empty and I do not toss in pebbles for fun.

Mark, from the Senate, comes to confiscate the magic pebbles, which are ordinary stones. My apprentice Autrey says the magic is in the level of pebbles. Mark fills his own bucket to a level he likes, and decides that ``magic'' is ``completely arbitrary and subjective.'' I ask what the pebbles do. They tell us when all the sheep are home.

An empty bucket is magical ``if and only if the pastures are empty,'' and I compare Tarski's sentence about snow. \nb{Rival theories of truth all accept Tarski's sentences, so the comparison separates the bucket from the pastures without favouring the correspondence view. The account rests on one kind of case, a count of visible objects checked the next morning. The first standard objection to correspondence theories is that they fit such cases and strain on others, such as moral and logical truths. The post does not take it up.}

Mark says sheep cannot interact with buckets, and proposes ``the coherentist theory of magic buckets'': magic comes from many buckets agreeing. He will not fill his family's buckets from ours, because they have always agreed with one another. \nb{That a false proposition can cohere with some set of propositions is Russell's objection to coherence theories, from 1907. The coherence theories described in the Stanford Encyclopedia take the relevant set to be, for example, what people actually believe, or would believe at the limit of inquiry. Neither is one family's collection closed to new buckets.} I show him a cloth on the gate that drops a pebble whenever a sheep brushes it. Mark asks at exactly which point the pebble becomes magic. I get confused, and answer that the point of the system is to keep track of sheep.

Autrey tries ``represent,'' ``special causal powers'' and ``about-ness,'' and I mock these words as magic under other names. I give Mark a pebble, drop in another from the ground, and declare Mark's hand part of the bucket. What a pebble stands for depends on its place in the count. \nb{Philosophers call accounts like this, on which a state is about what it has the job of tracking, teleological theories of content. The post names none of this work.} Adding a pebble does not make a sheep come home; I tested this with paperclips and dollar bills.

Mark says the trick failed because I did not believe in it. I answer: ``The bucket method works whether or not you believe in it.'' ``The correspondence between reality and my beliefs comes from reality controlling my beliefs, not the other way around.'' \nb{That is realism, the claim that the world exists apart from belief and causes it. It is not the same as the correspondence theory: deflationists, who reject that theory, are typically realists too.}

Mark suggests that calling a belief ``true'' only says you ``really really'' believe it. I am ``not entirely sure'' where reality comes from, but I need a name for whatever determines my results when even my best hypotheses are surprised. \nb{Peirce, who first defined and defended pragmatism, defined the real in 1878 as ``that whose characters are independent of what anybody may think them to be.'' Here the shepherd is close to Peirce. Mark's ``really really'' resembles the view of Richard Rorty, also a pragmatist.} Mark replies that whatever I say about reality is a belief, so only beliefs exist. I answer that whatever you eat is in your mouth, so there is no food, only mouths. \nb{David Stove made the same kind of parody, with oysters, of an argument of the same form.}

Inspector Darwin arrives. Mark now says there is no truth, that truth is a way to impose beliefs, that every culture has its own, and that this holds everywhere, and ``I insist that you agree.'' Autrey asks why he should believe Mark if nothing is true. Mark says his claims are ``fully general excuses'' that he applies ``selectively.'' The problem has a name, the ``reflexivity problem,'' and Mark gives no answer. \nb{Autrey's question goes back to Plato's reply to Protagoras. The literature Mark invokes does offer answers, and the question hits global relativism, which is much less common than the local kind. Since Mark admits he has no position, the dialogue's replies reach only Mark.}

Darwin proposes that Mark believe he can fly and step off the cliff. Autrey starts to object that everyone would then live in a private universe. Mark goes over the cliff, and Darwin lowers the frequency of Mark's alleles. Autrey and I go back to bringing in the sheep.
